\documentclass{a0poster}

% ------------------------------------------------------------------
%  Paquetes
% ------------------------------------------------------------------
\usepackage[utf8]{inputenc}
\usepackage[T1]{fontenc}
\usepackage{lmodern}                                  % fuentes escalables
\usepackage[spanish,es-nodecimaldot,es-tabla]{babel}  % punto decimal en modo matemático, "Tabla"

% Tamaño real del póster: 90 cm (ancho) x 130 cm (alto), vertical
% Para tamaño DIN A0 estándar: paperwidth=84.1cm, paperheight=118.9cm
% Para tamaño DIN A1: paperwidth=59.4cm, paperheight=84.1cm
\usepackage[
  paperwidth=90cm, paperheight=130cm,
  left=2cm, right=2cm,
  top=21 cm,    % cabecera (20.5 cm) + margen
  bottom=4cm     % pie de página
]{geometry}

\usepackage{amsmath, amssymb}
\usepackage{graphicx}
\usepackage{xcolor}
\usepackage{tikz}
\usepackage{pgfplots}
\usepackage{tcolorbox}
\usepackage{booktabs}
\usepackage{enumitem}
\usepackage{microtype}
\usepackage[font=large, labelfont={bf,color=UDLAPorange}, labelsep=period, skip=1.5ex]{caption}
\captionsetup[figure]{font=Large}

\pgfplotsset{compat=1.16}

% ------------------------------------------------------------------
%  Paleta de Colores (Personalizable)
% ------------------------------------------------------------------
% Colores predeterminados (estilo institucional UDLAP):
\definecolor{UDLAPorange}{RGB}{242,122,36} % Color primario / acentos
\definecolor{UDLAPdark}{RGB}{51,51,51}     % Color secundario / textos
\definecolor{UDLAPbg}{RGB}{253,240,230}   % Fondo tenue para cajas o realces

% Para cambiar a otra universidad o combinación de colores:
% UNAM: \definecolor{UDLAPorange}{RGB}{0,43,127}  % Azul UNAM
%       \definecolor{UDLAPdark}{RGB}{198,153,35}  % Oro UNAM
% IPN:  \definecolor{UDLAPorange}{RGB}{108,19,42}  % Guinda IPN

% ------------------------------------------------------------------
% TAMAÑO DE LETRA DEL CUERPO \normalsize (24.9 pt), \large (29.9 pt),
%  \Large (35.8 pt), \LARGE (38 pt), \huge (51.6 pt)
% ------------------------------------------------------------------
\newcommand{\bodysize}{\Large}

% ------------------------------------------------------------------
%  Ajustes generales
% ------------------------------------------------------------------
\pagestyle{empty}
\setlength{\parindent}{0pt}
\setlength{\parskip}{0.4\baselineskip}
\emergencystretch=3em
\setlength{\heavyrulewidth}{2pt}
\setlength{\lightrulewidth}{1pt}
\setlist{itemsep=0.4ex, topsep=0.4ex, parsep=0pt}

% Unidades abreviadas
\newcommand{\mm}{\,\mathrm{mm}}
\newcommand{\nm}{\,\mathrm{nm}}

% Ancho de cada columna
\newlength{\colw}
\setlength{\colw}{0.488\textwidth}

% ------------------------------------------------------------------
%  Ruta de figuras
% ------------------------------------------------------------------
\graphicspath{{figures/}{./}}

% Imagen con respaldo: busca primero en figures/ y luego en el directorio raíz.
% Si el archivo no existe, dibuja un recuadro de aviso en lugar de interrumpir la compilación.
\newcommand{\safegraphic}[2][]{%
  \IfFileExists{figures/#2}{\includegraphics[#1]{figures/#2}}{%
    \IfFileExists{#2}{\includegraphics[#1]{#2}}{%
      \fbox{\parbox[c][6cm][c]{12cm}{\centering\color{gray}\sffamily
        Falta el archivo:\\\texttt{\detokenize{#2}}}}}}}

% Tarjeta de sección: panel blanco con barra de título naranja
\newtcolorbox{card}[1]{%
  colback=white, colframe=UDLAPorange,
  coltitle=white, colbacktitle=UDLAPorange,
  title={#1}, fonttitle=\sffamily\bfseries\Huge, halign title=center,
  boxrule=1.5pt, arc=5mm,
  left=7mm, right=7mm, top=5mm, bottom=5mm,
  toptitle=3mm, bottomtitle=3mm,
  before skip=0pt, after skip=0.8cm}

% Datos de autoría
\newcommand{\posterauthors}{%
  Alan Ivan Ortega Contreras\textsuperscript{1}\quad
  Luz Stefanía Pérez Huesca\textsuperscript{1}\quad
  Ángel Gabriel Munguía Velázquez\textsuperscript{1}}

% ------------------------------------------------------------------
%  Cabecera y pie de página (se dibujan una sola vez, sobre la página)
% ------------------------------------------------------------------
\newcommand{\posterheader}{%
\begin{tikzpicture}[remember picture, overlay]
  % --- Franja blanca superior (8 cm) ---
  \fill[white] (current page.north west) rectangle ([yshift=-8cm]current page.north east);

  % --- Franja naranja con el título (12.5 cm) ---
  \fill[UDLAPorange] ([yshift=-8cm]current page.north west)
                     rectangle ([yshift=-20.5cm]current page.north east);

  % --- Logo UDLAP (izquierda) ---
  \node[anchor=west, inner sep=0pt] at ([xshift=3cm, yshift=-4cm]current page.north west)
    {\safegraphic[height=4.5cm]{udlap_logo.png}};

  % --- Institución (derecha) ---
  \node[anchor=east, align=right, text width=52cm, inner sep=0pt]
    at ([xshift=-3cm, yshift=-4cm]current page.north east) {%
    {\sffamily\bfseries\color{UDLAPorange}\huge Universidad de las Américas Puebla\par}
    \vspace{1ex}
    {\sffamily\bfseries\color{UDLAPdark}\LARGE Depto.\ de Actuaría, Física y Matemáticas\par}};

  % --- Título, autores y afiliación (sobre la franja naranja) ---
  \node[anchor=center, align=center, text=white, text width=84cm, inner sep=0pt]
    at ([yshift=-14.25cm]current page.north) {%
    \linespread{0.9}\selectfont
    {\sffamily\bfseries\Huge Interferencia con el espejo de Fresnel:\\[0pt]
     determinación de la longitud de onda de un láser He--Ne\par}
    \vspace{1.5ex}
    {\sffamily\LARGE \posterauthors\par}
    \vspace{0.5ex}
    {\sffamily\Large \textsuperscript{1}Depto.\ de Actuaría, Física y Matemáticas,
     Universidad de las Américas Puebla, Puebla, México 72810\par}};

  % --- Pie de página con curso y contacto ---
  \fill[UDLAPorange] ([yshift=0cm]current page.south west)
                     rectangle ([yshift=5cm]current page.south east);
  \node[anchor=west, align=left, text=white, inner sep=0pt]
    at ([xshift=3cm, yshift=2.5cm]current page.south west)
    {\linespread{0.9}\selectfont
     \sffamily\bfseries\LARGE Óptica Experimental, UDLAP\\[0pt]
     \large O26-LFA3032-1};
  \node[anchor=east, align=right, text=white, inner sep=0pt]
    at ([xshift=-3cm, yshift=2.5cm]current page.south east)
    {\linespread{0.9}\selectfont
     \sffamily\bfseries\Large Contacto\\[0pt]
     \large\mdseries alan.ortegacs@udlap.mx \quad luz.perezha@udlap.mx\\
     angel.munguiavz@udlap.mx};
\end{tikzpicture}%
}

% ==================================================================
\begin{document}

\pagecolor{UDLAPbg}
\posterheader

\bodysize

% ------------------------------------------------------------------
%  Abstract (ancho completo)
% ------------------------------------------------------------------
\begin{card}{Abstract}
Se estudió la interferencia de dos fuentes virtuales coherentes mediante el espejo de Fresnel. A partir de la separación entre franjas y de la separación geométrica de las fuentes se obtuvo una longitud de onda de \((670\pm70)\nm\), con un error del \(6.3\%\) respecto al valor nominal de \(632.8\nm\).

\textcolor{UDLAPorange}{\textbf{Keywords:}} interferencia, espejo de Fresnel, fuentes virtuales, longitud de onda, láser He--Ne.
\end{card}

\noindent
% ==================================================================
%  Columna izquierda
% ==================================================================
\begin{minipage}[t]{\colw}
\vspace{0pt}

\begin{card}{Objetivo e Introducción}
\textbf{Objetivo.} Determinar la longitud de onda de un láser He--Ne a partir de la separación entre franjas de interferencia (\(\bar d\)) y de la separación proyectada de las fuentes virtuales (\(A\)).

\textbf{Introducción.} El espejo de Fresnel genera dos fuentes virtuales , \(S_1'\) y \(S_2'\), mediante la reflexión de una fuente puntual. Esto permite observar un patrón de interferencia estable y evidenciar la naturaleza ondulatoria de la luz.
\end{card}

\begin{card}{Marco teórico}
\begin{minipage}[t]{0.40\linewidth}
\vspace{0pt}
El haz reflejado por los dos semi-espejos ligeramente inclinados genera dos ondas que se originan en las fuentes virtuales \(S_1'\) y \(S_2'\). Los máximos de intensidad cumplen \(\Delta s=n\lambda\), con \(\Delta s\) la diferencia de trayectoria óptica. Para una pantalla a gran distancia \(L\) (Fig.~\ref{fig:diagrama}), con \(a\) la separación entre las fuentes y \(\bar d\) la distancia medida entre franjas,
\begin{equation}
\lambda=a\,\frac{\bar d}{L}.
\label{eq:lambda}
\end{equation}
\end{minipage}\hfill
\begin{minipage}[t]{0.58\linewidth}
\vspace{0pt}
\centering
\safegraphic[width=\linewidth,height=18cm,keepaspectratio]{Diagrama_Fresnel.png}
\captionof{figure}{Diagrama de rayos del espejo de Fresnel.}
\label{fig:diagrama}
\end{minipage}

\begin{minipage}[t]{0.40\linewidth}
\vspace{0pt}
Para determinar \(a\), se proyectan las fuentes con una segunda lente (Fig.~\ref{fig:trayectoria}); por triángulos semejantes,
\begin{equation}
a=A\,\frac{L_1}{L_2},
\label{eq:a}
\end{equation}
con \(A\) la separación entre las imágenes proyectadas, \(L_1\) la distancia de las fuentes a la lente y \(L_2\) la de la lente a la pantalla. Combinando ambas ecuaciones se obtiene un método puramente geométrico para estimar \(\lambda\).
\end{minipage}\hfill
\begin{minipage}[t]{0.58\linewidth}
\vspace{0pt}
\centering
\vspace{-5cm}
\safegraphic[width=\linewidth,height=24cm,keepaspectratio]{trayectoria.png}
\captionof{figure}{Trayectoria del haz y proyección con la lente.}
\label{fig:trayectoria}
\end{minipage}

\end{card}

\begin{card}{Montaje experimental}
\begin{center}
\safegraphic[width=0.72\linewidth,height=26cm,keepaspectratio]{PHOTO-2026-10-02-07-00-35.jpg}
\captionof{figure}{Montaje: láser He--Ne, lentes, espejo de Fresnel y pantalla traslúcida.}
\label{fig:montaje}
\end{center}
\end{card}

\end{minipage}\hfill
% ==================================================================
%  Columna derecha
% ==================================================================
\begin{minipage}[t]{\colw}
\vspace{0pt}

\begin{card}{Resultados y Análisis}
\begin{minipage}[t]{0.48\linewidth}
\vspace{0pt}
\centering
\safegraphic[width=\linewidth,height=11cm,keepaspectratio]{IMG_5736.jpeg}
\captionof{figure}{Fuentes virtuales proyectadas.}
\label{fig:puntos}
\end{minipage}\hfill
\begin{minipage}[t]{0.48\linewidth}
\vspace{0pt}
\centering
\safegraphic[width=\linewidth,height=12.5cm,keepaspectratio]{IMG_5734.jpeg}
\captionof{figure}{Patrón de interferencia.}
\label{fig:patron}
\end{minipage}
\par\vspace{1ex}
Se midieron cinco separaciones \(d_k\) entre máximos consecutivos (Fig.~\ref{fig:grafica_d}) y las distancias del montaje (Tabla~\ref{tab:distancias}).

\begin{center}
\begin{tikzpicture}
\begin{axis}[
  width=0.92\linewidth, height=11.5cm,
  xlabel={Medición \(k\)}, ylabel={\(d_k\) (mm)},
  xmin=0.5, xmax=5.5, ymin=0.2, ymax=0.6,
  xtick={1,2,3,4,5}, ytick={0.2,0.3,0.4,0.5,0.6},
  grid=major, grid style={gray!25},
  label style={font=\Large}, tick label style={font=\large},
  legend style={font=\large, draw=none, fill=white, fill opacity=0.85, text opacity=1},
  legend pos=north east]
  \fill[UDLAPorange!25] (axis cs:0.5,0.355) rectangle (axis cs:5.5,0.429);
  \addplot+[only marks, mark=*, mark size=4pt, color=UDLAPdark,
            nodes near coords={\pgfmathprintnumber[fixed,fixed zerofill,precision=2]{\pgfplotspointmeta}},
            every node near coord/.append style={font=\large, anchor=south, yshift=4pt},
            error bars/.cd, y dir=both, y explicit]
    coordinates {(1,0.36) +- (0,0.01) (2,0.46) +- (0,0.01) (3,0.50) +- (0,0.01)
                 (4,0.32) +- (0,0.01) (5,0.32) +- (0,0.01)};
  \addplot[dashed, very thick, UDLAPorange] coordinates {(0.5,0.392) (5.5,0.392)};
  \legend{\(d_k\), \(\bar d\)}
\end{axis}
\end{tikzpicture}
\captionof{figure}{Separaciones medidas (\(\pm0.01\mm\)), promedio \(\bar d\) y banda \(\pm\sigma_{\bar d}\).}
\label{fig:grafica_d}
\end{center}

\vspace{-1ex}
\[
\bar d=\frac{1}{N}\sum_{k=1}^{N}d_k=0.392\mm,\qquad
\sigma_{\bar d}=\frac{\sigma_d}{\sqrt N}=0.037\mm,
\]
con \(N=5\) y \(\sigma_d\) la desviación estándar de las \(d_k\); es decir, \mbox{\(\bar d=(0.39\pm0.04)\mm\)}.

\begin{center}
\Large
\setlength{\tabcolsep}{10pt}
\renewcommand{\arraystretch}{1.3}
\begin{tabular}{ccccc}
\toprule
\(A\) & \(L_0\) & \(L_2\) & \(L_1\) & \(L\)\\
\midrule
\(5.0\pm0.1\) & \(290.0\pm0.5\) & \(780.0\pm0.5\) & \(285.0\pm0.5\) & \(1065.0\pm0.7\)\\
\bottomrule
\end{tabular}
\captionof{table}{Distancias del montaje (mm). \(L_0\): entre la lente de \(5\mm\) y la de proyección; \(L_1=L_0-5\mm\) (la fuente puntual está en el foco de la lente de \(5\mm\)); \(L=L_1+L_2\).}
\label{tab:distancias}
\end{center}

Con los valores sin redondear:
\begin{align*}
a&=A\,\frac{L_1}{L_2}=5.0\mm\cdot\frac{285.0\mm}{780.0\mm}=1.8269\mm,\\
\lambda&=a\,\frac{\bar d}{L}=1.8269\mm\cdot\frac{0.392\mm}{1065.0\mm}=672.4\nm.
\end{align*}

Propagando las incertidumbres (variables independientes),
\[
\delta\lambda=\lambda\sqrt{\sum_i\left(\frac{\delta x_i}{x_i}\right)^{2}},
\qquad x_i\in\{A,L_1,L_2,L,\bar d\}.
\]
Domina \(\bar d\) (\(\approx9.5\%\)), seguida de \(A\) (\(2\%\)): \(\delta\lambda\approx65\nm\approx70\nm\).

\begin{tcolorbox}[colback=UDLAPorange!10, colframe=UDLAPorange, boxrule=2pt, arc=3mm]
\[
\lambda=(670\pm70)\nm,\qquad
\%\mathrm{error}=\frac{|672.4-632.8|}{632.8}\times100\%=6.3\%
\]
\end{tcolorbox}
\end{card}

\begin{card}{Conclusiones}
El espejo de Fresnel generó fuentes virtuales coherentes separadas \(a\approx1.83\mm\). Se obtuvo \(\lambda=(670\pm70)\nm\), con un error de \(6.3\%\) respecto al valor nominal del láser He--Ne, lo que valida cuantitativamente la relación entre el espaciamiento de las franjas y la longitud de onda.
\end{card}

\begin{card}{Referencias}
\normalsize
\begin{enumerate}[label={[\arabic*]}, leftmargin=*, labelsep=1ex, itemsep=1ex]
\item LD DIDACTIC GmbH, \emph{Interference at a Fresnel's mirror with an He--Ne laser}, Physics Leaflets P5.3.2.1, pp.~1--3.
\item LEYBOLD DIDACTIC GmbH, \emph{He--Ne laser, linearly polarised (471 830)}, Instruction sheet, 06/04-W97-Kem, pp.~1--2.
\end{enumerate}
\end{card}

\end{minipage}

\end{document}